\section{动态组织选择、Bellman 系统与值函数存在唯一性}

本阶段研究 benchmark 模型在固定价格下的动态模块。全文在这一阶段内把 $(w,p_m,M)$ 视为给定。因此，这里证明的是固定价格 Bellman 系统的存在性、唯一性、策略存在性以及数值可解性；这里还不是关于完整 stationary competitive equilibrium 的存在性证明。

\subsection{Benchmark Fidelity Check}

进入本阶段 Bellman 系统的精确静态收益对象是
\[
\Pi_A(z;w,M),\quad \Pi_B(z;p_m,M),\quad \Pi_C(z;p_m,w,M),
\]
它们必须保持这些精确的 argument list，不能改写为 $\Pi_j(z;w,p_m,M)$。本阶段要解出的动态内生对象是
\[
V_A(z;M),\quad V_B(z;M),\quad V_C(z;M),
\]
以及由它们派生的 choice-specific values $W_{sj}(z;M)$ 和 choice regions。

这里需要明确指出一处 benchmark source 之间的记号张力。主稿动态部分有时使用通用写法 $\Pi_j(z;w,p_m,M)$，但 canonical 的静态对象始终是
\[
\Pi_A(z;w,M),\qquad \Pi_B(z;p_m,M),\qquad \Pi_C(z;p_m,w,M).
\]
这两种写法在字面上并不完全相同。本阶段为保持 benchmark fidelity，继续保留对象专属的 payoff 定义，并把 $\Pi_j(z;w,p_m,M)$ 仅解释为“对应类型 $j$ 的优化后静态收益”的简写，而不是对各对象 argument list 的重新定义。在这个澄清之下，本阶段的固定价格 Bellman 推导可以忠实完成。

\subsection{Part 1：动态问题、时序与 choice-specific value}

固定一个 incumbent firm，它在期初继承的组织类型为 $s\in\{A,B,C\}$，继承的 capability state 为 $z\in Z$。该企业在当前期的可行组织行动集合是
\[
\mathcal J(s)=\{A,B,C,\mathrm{exit}\}.
\]
为后文简写，记
\[
\mathcal A:=\{A,B,C,\mathrm{exit}\}.
\]
benchmark 的时序是：
\begin{enumerate}
    \item 企业先选择当期组织行动 $j\in\mathcal J(s)$。
    \item 如果 $j\neq s$，企业立即支付切换成本 $\kappa_{sj}(M)$。如果 $j=s$，则没有组织重构，因此把 $\kappa_{ss}(M)$ 解释为 $0$。
    \item 如果企业选择活跃类型 $j\in\{A,B,C\}$，它随后解决该类型对应的当期经营问题。Phase 2 已经把这些期内问题解完，并分别汇总为优化后的静态收益对象 $\Pi_A,\Pi_B,\Pi_C$。
    \item 如果企业选择 exit，则它获得被标准化为 $0$ 的退出价值，并且在这个行业内的动态问题中不再有 continuation value。
    \item 如果企业选择活跃类型 $j\in\{A,B,C\}$，则下一期 capability state $z'$ 按照类型专属转移核 $Q_j(\cdot\mid z,M)$ 抽取，下一期以继承组织类型 $j$ 开始。
\end{enumerate}

为什么组织选择必须先于当期经营控制？原因不是写法偏好，而是 benchmark 结构本身要求如此。组织形式决定了本期究竟面对哪一个静态问题：若选择 $A$，就面对 type-$A$ 的期内控制；若选择 $B$，就面对 type-$B$ 的期内控制；若选择 $C$，就面对 type-$C$ 的 manufacturing 决策；若选择 exit，则根本没有经营控制。因此，组织形式是外层动态选择，经营控制是内层期内选择，而内层问题已经在 Phase 2 中被压缩进 $\Pi_A,\Pi_B,\Pi_C$。

对继承类型 $s\in\{A,B,C\}$，定义 benchmark 的 choice-specific values：
\begin{align*}
W_{sA}(z;M)
&=
\Pi_A(z;w,M)-\kappa_{sA}(M)
+\beta\int_Z V_A(z';M)\,Q_A(dz'\mid z,M),\\
W_{sB}(z;M)
&=
\Pi_B(z;p_m,M)-\kappa_{sB}(M)
+\beta\int_Z V_B(z';M)\,Q_B(dz'\mid z,M),\\
W_{sC}(z;M)
&=
\Pi_C(z;p_m,w,M)-\kappa_{sC}(M)
+\beta\int_Z V_C(z';M)\,Q_C(dz'\mid z,M),\\
W_{s,\mathrm{exit}}(z;M)
&=0.
\end{align*}
这里每一个 $W_{sj}(z;M)$ 都包含四项：选择组织形式 $j$ 后的当期优化收益、当前由类型 $s$ 切换到类型 $j$ 的成本、折现因子 $\beta$，以及下一期以类型 $j$ 进入时的预期 continuation value。

\subsection{Part 2：Bellman 方程}

对每个继承类型 $s\in\{A,B,C\}$，Bellman 问题是
\[
V_s(z;M)=\max_{j\in\mathcal J(s)}W_{sj}(z;M).
\]
因为每个当前类型的可行集合都是 $\mathcal J(s)=\{A,B,C,\mathrm{exit}\}$，所以 benchmark Bellman 系统写为
\begin{align*}
V_A(z;M)&=\max\{W_{AA}(z;M),W_{AB}(z;M),W_{AC}(z;M),0\},\\
V_B(z;M)&=\max\{W_{BA}(z;M),W_{BB}(z;M),W_{BC}(z;M),0\},\\
V_C(z;M)&=\max\{W_{CA}(z;M),W_{CB}(z;M),W_{CC}(z;M),0\}.
\end{align*}
把上面各个 $W_{sj}$ 的定义逐项代回去，就得到完全展开的系统：
\begin{align*}
V_A(z;M)
&=
\max\Bigl\{
\Pi_A(z;w,M)-\kappa_{AA}(M)+\beta\!\int_Z\!V_A(z';M)Q_A(dz'\mid z,M),\\
&\hspace{1.75cm}
\Pi_B(z;p_m,M)-\kappa_{AB}(M)+\beta\!\int_Z\!V_B(z';M)Q_B(dz'\mid z,M),\\
&\hspace{1.75cm}
\Pi_C(z;p_m,w,M)-\kappa_{AC}(M)+\beta\!\int_Z\!V_C(z';M)Q_C(dz'\mid z,M),\,0
\Bigr\},\\
V_B(z;M)
&=
\max\Bigl\{
\Pi_A(z;w,M)-\kappa_{BA}(M)+\beta\!\int_Z\!V_A(z';M)Q_A(dz'\mid z,M),\\
&\hspace{1.75cm}
\Pi_B(z;p_m,M)-\kappa_{BB}(M)+\beta\!\int_Z\!V_B(z';M)Q_B(dz'\mid z,M),\\
&\hspace{1.75cm}
\Pi_C(z;p_m,w,M)-\kappa_{BC}(M)+\beta\!\int_Z\!V_C(z';M)Q_C(dz'\mid z,M),\,0
\Bigr\},\\
V_C(z;M)
&=
\max\Bigl\{
\Pi_A(z;w,M)-\kappa_{CA}(M)+\beta\!\int_Z\!V_A(z';M)Q_A(dz'\mid z,M),\\
&\hspace{1.75cm}
\Pi_B(z;p_m,M)-\kappa_{CB}(M)+\beta\!\int_Z\!V_B(z';M)Q_B(dz'\mid z,M),\\
&\hspace{1.75cm}
\Pi_C(z;p_m,w,M)-\kappa_{CC}(M)+\beta\!\int_Z\!V_C(z';M)Q_C(dz'\mid z,M),\,0
\Bigr\}.
\end{align*}
这个问题之所以是 recursive 的，是因为我们今天要求解的未知对象，也就是 value function，本身又出现在下一期 continuation term 中。

\subsection{Part 3：函数空间、范数与 Bellman 算子}

由 Phase 1，状态空间 $Z$ 是非空紧致 metric space，$\mathcal B(Z)$ 是它的 Borel sigma 代数。定义
\[
\mathcal C_b(Z)=\{f:Z\to\mathbb R\mid f\text{ 有界且连续}\},
\qquad
\|f\|_\infty=\sup_{z\in Z}|f(z)|.
\]
由于 benchmark 有三个活跃组织类型，所以候选值函数空间是
\[
\mathcal V:=\mathcal C_b(Z)^3.
\]
对任意候选三元组 $U=(U_A,U_B,U_C)\in\mathcal V$，定义乘积 sup 范数
\[
\|U\|:=\max\{\|U_A\|_\infty,\|U_B\|_\infty,\|U_C\|_\infty\}.
\]

为了紧凑地定义一个 Bellman 算子，我们引入只用于证明的辅助记号。对 $j\in\{A,B,C\}$，定义
\[
(\mathcal Q_jU)(z):=\int_Z U_j(z')Q_j(dz'\mid z,M),
\qquad
(\mathcal Q_{\mathrm{exit}}U)(z):=0.
\]
对继承类型 $s\in\{A,B,C\}$ 与动作 $a\in\{A,B,C,\mathrm{exit}\}$，定义
\[
\mathcal W_{sa}(z;M;U):=
\begin{cases}
\Pi_A(z;w,M)-\kappa_{sA}(M)+\beta(\mathcal Q_AU)(z), & a=A,\\
\Pi_B(z;p_m,M)-\kappa_{sB}(M)+\beta(\mathcal Q_BU)(z), & a=B,\\
\Pi_C(z;p_m,w,M)-\kappa_{sC}(M)+\beta(\mathcal Q_CU)(z), & a=C,\\
0, & a=\mathrm{exit}.
\end{cases}
\]
在固定点 $V^*$ 处，benchmark 的 choice-specific values 由
\[
W_{sa}(z;M)=\mathcal W_{sa}(z;M;V^*)
\]
恢复出来。于是 Bellman 算子定义为
\[
(\mathcal T_MU)_s(z):=\max_{a\in\{A,B,C,\mathrm{exit}\}}\mathcal W_{sa}(z;M;U),
\qquad s\in\{A,B,C\}.
\]
因此，固定价格 Bellman 问题等价于固定点方程 $V=\mathcal T_MV$。

\subsection{假设}

\begin{assumption}[R1: 流收益有界连续]
优化后的静态收益满足 $\Pi_A(\cdot;w,M)\in\mathcal C_b(Z)$，$\Pi_B(\cdot;p_m,M)\in\mathcal C_b(Z)$，且 $\Pi_C(\cdot;p_m,w,M)\in\mathcal C_b(Z)$。
\end{assumption}

\begin{assumption}[R2: Feller 概率转移核]
对每个 $j\in\{A,B,C\}$ 和每个 $(z,M)$，$Q_j(\cdot\mid z,M)$ 都是定义在 $Z$ 上的 Borel 概率测度。此外，对每个 $f\in\mathcal C_b(Z)$，映射
\[
z\mapsto\int_Z f(z')Q_j(dz'\mid z,M)
\]
在 $Z$ 上连续。
\end{assumption}

\begin{assumption}[R3: 折现]
$\beta\in(0,1)$。
\end{assumption}

\begin{assumption}[R4: 切换成本有限]
对所有 $s,j\in\{A,B,C\}$，$|\kappa_{sj}(M)|<\infty$。
\end{assumption}

\subsection{Part 4：Bellman 系统的 well-posedness}

\begin{lemma}
在 R1--R4 下，$(\mathcal C_b(Z),\|\cdot\|_\infty)$ 是完备的。因此，$(\mathcal V,\|\cdot\|)$ 也是完备的。
\end{lemma}

\begin{proof}
取 $(\mathcal C_b(Z),\|\cdot\|_\infty)$ 中任意一个 Cauchy 序列 $\{f^n\}$。对每个 $\varepsilon>0$，存在 $N(\varepsilon)$，使得
\[
\|f^n-f^m\|_\infty<\varepsilon
\qquad\text{对所有 }n,m\ge N(\varepsilon).
\]
固定任意 $z\in Z$，则
\[
|f^n(z)-f^m(z)|\le \|f^n-f^m\|_\infty<\varepsilon
\qquad\text{对所有 }n,m\ge N(\varepsilon).
\]
因此 $\{f^n(z)\}$ 在 $\mathbb R$ 中是 Cauchy 序列，从而收敛到某个极限 $f(z)\in\mathbb R$。

下面证明收敛是一致的。固定 $n\ge N(\varepsilon)$。对所有 $m\ge N(\varepsilon)$，
\[
\sup_{z\in Z}|f^n(z)-f^m(z)|<\varepsilon.
\]
令 $m\to\infty$，由于 $f^m(z)\to f(z)$ 逐点成立，所以
\[
|f^n(z)-f(z)|\le \varepsilon
\qquad\text{对每个 }z\in Z.
\]
对 $z$ 取上确界，得到
\[
\|f^n-f\|_\infty\le \varepsilon.
\]
所以 $f^n\to f$ 在 $Z$ 上一致收敛。

由于每个 $f^n$ 都连续，且收敛是一致的，所以根据 Uniform Limit Theorem，极限函数 $f$ 仍然连续。又因为该序列在 sup 范数下是 Cauchy，存在 $N$ 使得当 $n\ge N$ 时，$\|f^n-f^N\|_\infty\le 1$。令 $n\to\infty$，得到 $\|f-f^N\|_\infty\le 1$，因此
\[
\|f\|_\infty\le \|f-f^N\|_\infty+\|f^N\|_\infty\le 1+\|f^N\|_\infty<\infty.
\]
于是 $f\in\mathcal C_b(Z)$，说明 $\mathcal C_b(Z)$ 完备。

对于乘积空间，令 $U^n=(U_A^n,U_B^n,U_C^n)$ 在 $\mathcal V$ 中是 Cauchy。则每个坐标序列在 $\mathcal C_b(Z)$ 中也是 Cauchy，因为
\[
\|U_A^n-U_A^m\|_\infty\le \|U^n-U^m\|,
\quad
\|U_B^n-U_B^m\|_\infty\le \|U^n-U^m\|,
\quad
\|U_C^n-U_C^m\|_\infty\le \|U^n-U^m\|.
\]
每个坐标都在 $\mathcal C_b(Z)$ 中收敛，因此三者合起来给出 $\mathcal V$ 中的极限。所以 $(\mathcal V,\|\cdot\|)$ 完备。
\end{proof}

\begin{lemma}
在 R1--R4 下，对任意 $U\in\mathcal V$，都有 $\mathcal T_MU\in\mathcal V$。
\end{lemma}

\begin{proof}
固定 $U\in\mathcal V$ 以及继承类型 $s\in\{A,B,C\}$。

对每个活跃动作 $j\in\{A,B,C\}$，
\begin{align*}
|(\mathcal Q_jU)(z)|
&=
\left|\int_Z U_j(z')Q_j(dz'\mid z,M)\right|\\
&\le
\int_Z |U_j(z')|Q_j(dz'\mid z,M)\\
&\le
\int_Z \|U_j\|_\infty Q_j(dz'\mid z,M)\\
&=
\|U_j\|_\infty.
\end{align*}
所以每个 continuation term 都有界。由 R2，且因为 $U_j\in\mathcal C_b(Z)$，每个 continuation term 也是连续的。因此，对活跃动作有
\[
\mathcal W_{sj}(\cdot;M;U)\in\mathcal C_b(Z).
\]
对 exit 动作，$\mathcal W_{s,\mathrm{exit}}(\cdot;M;U)\equiv 0$，它同样有界且连续。

由于 $(\mathcal T_MU)_s$ 是有限个有界连续函数的最大值，所以 $(\mathcal T_MU)_s\in\mathcal C_b(Z)$。这对 $s=A,B,C$ 都成立，因此 $\mathcal T_MU\in\mathcal V$。
\end{proof}

\begin{lemma}
在 R1--R4 下，对所有 $U,\widehat U\in\mathcal V$，
\[
\|\mathcal T_MU-\mathcal T_M\widehat U\|\le\beta\|U-\widehat U\|.
\]
\end{lemma}

\begin{proof}
固定继承类型 $s$ 与状态 $z$。先证明基础不等式
\[
|\max_i x_i-\max_i y_i|\le\max_i|x_i-y_i|.
\]
令 $i_x$ 取到 $\max_i x_i$。则
\[
\max_i x_i-\max_i y_i=x_{i_x}-\max_i y_i\le x_{i_x}-y_{i_x}\le \max_i|x_i-y_i|.
\]
交换 $x$ 与 $y$ 的角色，就得到
\[
\max_i y_i-\max_i x_i\le \max_i|x_i-y_i|.
\]
合并两式，即得绝对值版本的不等式。

于是
\[
|(\mathcal T_MU)_s(z)-(\mathcal T_M\widehat U)_s(z)|
\le\max_{a\in\{A,B,C,\mathrm{exit}\}}|\mathcal W_{sa}(z;M;U)-\mathcal W_{sa}(z;M;\widehat U)|.
\]
对 exit 动作，差值为零。对活跃动作 $a=j\in\{A,B,C\}$，
\begin{align*}
|\mathcal W_{sj}(z;M;U)-\mathcal W_{sj}(z;M;\widehat U)|
&=\beta\left|\int_Z \bigl(U_j(z')-\widehat U_j(z')\bigr)Q_j(dz'\mid z,M)\right|\\
&\le\beta\int_Z |U_j(z')-\widehat U_j(z')|Q_j(dz'\mid z,M)\\
&\le\beta\|U_j-\widehat U_j\|_\infty\\
&\le\beta\|U-\widehat U\|.
\end{align*}
对动作取最大，对状态取上确界，再对继承类型取最大，就得到压缩不等式。
\end{proof}

\begin{theorem}
在 R1--R4 下，存在唯一的值函数三元组
\[
V^*=(V_A^*,V_B^*,V_C^*)\in\mathcal V
\]
使得
\[
\mathcal T_MV^*=V^*.
\]
等价地说，固定价格 Bellman 系统在 $\mathcal V$ 中有唯一解。
\end{theorem}

\begin{proof}
这里调用 Banach Fixed Point Theorem。该定理表述为：如果 $(X,d)$ 是完备度量空间，且 $T:X\to X$ 满足
\[
d(Tx,Ty)\le q\,d(x,y)\qquad\text{对所有 }x,y\in X
\]
其中 $q\in[0,1)$，那么 $T$ 在 $X$ 上存在唯一固定点，而且从任意初值出发的迭代都会收敛到该固定点。

现在逐条核对定理假设。第一，前面的完备性引理已经证明 $(\mathcal V,\|\cdot\|)$ 是完备度量空间。第二，上一条自映射引理已经证明 $\mathcal T_M:\mathcal V\to\mathcal V$。第三，压缩引理已经证明
\[
\|\mathcal T_MU-\mathcal T_M\widehat U\|\le \beta\|U-\widehat U\|
\]
对所有 $U,\widehat U\in\mathcal V$ 成立，而 R3 给出 $\beta\in(0,1)$，因此 $\beta<1$。

于是 Banach 定理的所有假设都满足。该定理推出存在唯一固定点 $V^*\in\mathcal V$，满足 $\mathcal T_MV^*=V^*$。由于固定点方程正是 Bellman 系统，所以值函数三元组唯一确定。
\end{proof}

\subsection{Part 5：策略存在性、set-valued policy 与 measurable selector}

对每个继承类型 $s\in\{A,B,C\}$，定义 benchmark 的 argmax correspondence：
\[
\Gamma_s(z;M):=\arg\max_{a\in\mathcal A}\mathcal W_{sa}(z;M;V^*).
\]

\begin{proposition}
对每个继承类型 $s\in\{A,B,C\}$：
\begin{enumerate}
    \item 对每个 $z\in Z$，$\Gamma_s(z;M)$ 都非空；
    \item $\Gamma_s(z;M)$ 可能是单值，也可能是多值，这取决于最优动作是否唯一；
    \item 对每个动作 $j\in\{A,B,C,\mathrm{exit}\}$，benchmark 的 choice region
    \[
    \mathcal Z_{sj}(M):=\{z\in Z:j\in\Gamma_s(z;M)\}
    \]
    是 Borel 可测的。
\end{enumerate}
\end{proposition}

\begin{proof}
因为动作集合 $\mathcal A$ 是有限集，而实值函数在有限集上总能达到最大值，所以 $\Gamma_s(z;M)$ 非空。如果有一个动作严格优于其余所有动作，那么该 correspondence 是单值的；如果有两个或以上动作在最大值处并列，则它是多值的。benchmark 在这一层并没有额外强制 deterministic tie-breaking。

下面证明 choice region 的可测性。固定 $j\in\{A,B,C,\mathrm{exit}\}$。根据定义，
\[
\mathcal Z_{sj}(M)
=
\bigcap_{\ell\in\mathcal A}
\{z\in Z:\mathcal W_{sj}(z;M;V^*)\ge \mathcal W_{s\ell}(z;M;V^*)\}.
\]
对每个竞争动作 $\ell$，差值函数
\[
z\mapsto \mathcal W_{sj}(z;M;V^*)-\mathcal W_{s\ell}(z;M;V^*)
\]
是连续的，因此集合
\[
\{z\in Z:\mathcal W_{sj}(z;M;V^*)-\mathcal W_{s\ell}(z;M;V^*)\ge 0\}
\]
是连续映射下闭集 $[0,\infty)$ 的原像，所以它是闭集。有限个闭集的交仍是闭集，所以 $\mathcal Z_{sj}(M)$ 是 Borel 可测的。
\end{proof}

\subsection{选择区域与 Phase 4 接口}

为与 benchmark 的 transition interface 保持一致，直接由 Bellman 问题定义 choice regions：
\[
\mathcal Z_{sj}(M)=\{z\in Z:j\in\Gamma_s(z;M)\},
\qquad s\in\{A,B,C\},\ j\in\{A,B,C,\mathrm{exit}\}.
\]
给定稳态测度 $\mu_s$，Phase 4 使用
\[
P_{sj}(M)=\frac{\mu_s(\mathcal Z_{sj}(M))}{\mu_s(Z)}.
\]
由于这些 choice regions 可测，只要 $\mu_s(Z)>0$，上式就有良好定义。

如果后续计算必须使用单值 deterministic policy，那么只能在显式声明 tie-breaking 约定之后，额外构造一个 selector。比如，在次序 $A\succ B\succ C\succ \mathrm{exit}$ 下，定义
\begin{align*}
D_{sA}(M)&:=\mathcal Z_{sA}(M),\\
D_{sB}(M)&:=\mathcal Z_{sB}(M)\setminus D_{sA}(M),\\
D_{sC}(M)&:=\mathcal Z_{sC}(M)\setminus\bigl(D_{sA}(M)\cup D_{sB}(M)\bigr),\\
D_{s,\mathrm{exit}}(M)&:=\mathcal Z_{s,\mathrm{exit}}(M)\setminus\bigl(D_{sA}(M)\cup D_{sB}(M)\cup D_{sC}(M)\bigr).
\end{align*}
然后令
\[
g_s(z;M)=
\begin{cases}
A, & z\in D_{sA}(M),\\
B, & z\in D_{sB}(M),\\
C, & z\in D_{sC}(M),\\
\mathrm{exit}, & z\in D_{s,\mathrm{exit}}(M)
\end{cases}
\]
就得到一个满足 $g_s(z;M)\in\Gamma_s(z;M)$ 的 Borel 可测 selector。这个 selector 只是辅助实现对象；benchmark-faithful 的策略对象仍然是 correspondence $\Gamma_s$，或者等价地说，是 measurable choice regions $\mathcal Z_{sj}(M)$。

\subsection{Bellman 与 value function 的求解过程}

前面的各个小节已经分别给出了 Bellman 方程、Bellman 算子、存在唯一性条件以及策略对象的定义。本小节把这些结果整理为一个连续的求解顺序，目的是明确说明：在固定价格 $(w,p_m,M)$ 下，如何从优化后静态收益对象出发，得到唯一的 value-function 三元组，并据此恢复后续阶段所需的策略与转移对象。

\paragraph{第一步：构造 choice-specific values。}
对继承类型 $s\in\{A,B,C\}$ 和状态 $z\in Z$，企业面对四个动作：
\[
\mathcal A=\{A,B,C,\mathrm{exit}\}.
\]
如果企业选择组织形式 $A$，那么它的当前总价值是
\[
W_{sA}(z;M)
=
\Pi_A(z;w,M)-\kappa_{sA}(M)+\beta\int_Z V_A(z';M)Q_A(dz'\mid z,M).
\]
如果企业选择 $B$ 或 $C$，定义完全类似：
\[
W_{sB}(z;M)
=
\Pi_B(z;p_m,M)-\kappa_{sB}(M)+\beta\int_Z V_B(z';M)Q_B(dz'\mid z,M),
\]
\[
W_{sC}(z;M)
=
\Pi_C(z;p_m,w,M)-\kappa_{sC}(M)+\beta\int_Z V_C(z';M)Q_C(dz'\mid z,M).
\]
如果企业选择退出，则
\[
W_{s,\mathrm{exit}}(z;M)=0.
\]
因此，固定价格动态问题的第一步不是直接猜测 $V_s$ 的 closed form，而是先把每一个当前可行动作对应的整期价值写清楚。这样处理之后，Bellman 问题的比较对象就不再是抽象的“组织选择”，而是四个明确定义的 choice-specific values。

\paragraph{第二步：把 value function 写成对 choice-specific values 的上包络。}
一旦四个 choice-specific values 被写出来，继承类型 $s$ 的 value function 就由
\[
V_s(z;M)=\max_{j\in\mathcal A}W_{sj}(z;M)
\]
给出。也就是说，
\[
V_A(z;M)=\max\{W_{AA}(z;M),W_{AB}(z;M),W_{AC}(z;M),0\},
\]
\[
V_B(z;M)=\max\{W_{BA}(z;M),W_{BB}(z;M),W_{BC}(z;M),0\},
\]
\[
V_C(z;M)=\max\{W_{CA}(z;M),W_{CB}(z;M),W_{CC}(z;M),0\}.
\]
因此，在这个 benchmark 中，“解 Bellman 系统”与“求出三个 value functions $V_A,V_B,V_C$”本质上是同一件事。因为一旦这三个 value functions 被求出，choice-specific values 的排序、argmax correspondences 以及 choice regions 都随之确定。

\paragraph{第三步：把 Bellman 系统改写为固定点问题。}
上面的三条方程之所以难，是因为当前期未知的 value function 又出现在 continuation term 里。例如，
\[
W_{sA}(z;M)
=
\Pi_A(z;w,M)-\kappa_{sA}(M)+\beta\int_Z V_A(z';M)Q_A(dz'\mid z,M)
\]
右边仍然含有未知函数 $V_A$。所以这里不是在解一个有限维代数方程组，而是在解一个函数方程组。更准确地说，我们是在寻找一个函数三元组
\[
V=(V_A,V_B,V_C)
\]
使它满足
\[
V=\mathcal T_MV.
\]
因此，“解 Bellman”在数学上等价于“求 Bellman 算子的固定点”。这一步把一个彼此递归的函数方程组转化成了一个定义在函数空间上的统一固定点问题。

\paragraph{第四步：用 Banach 定理得到唯一的 value function。}
为了证明这个 fixed point 真正存在，而且只有一个，我们把候选值函数放进函数空间
\[
\mathcal V=\mathcal C_b(Z)^3
\]
中，并在这个空间上定义 product sup norm。接着证明三件事：
\begin{enumerate}
    \item $(\mathcal V,\|\cdot\|)$ 是完备的；
    \item Bellman 算子 $\mathcal T_M$ 把 $\mathcal V$ 映回自身；
    \item Bellman 算子满足压缩不等式
    \[
    \|\mathcal T_MU-\mathcal T_M\widehat U\|\le \beta\|U-\widehat U\|.
    \]
\end{enumerate}
因为 $\beta\in(0,1)$，所以 $\mathcal T_M$ 是压缩映射。根据 Banach Fixed Point Theorem，存在唯一的固定点
\[
V^*=(V_A^*,V_B^*,V_C^*)\in\mathcal V
\]
满足
\[
\mathcal T_MV^*=V^*.
\]
这一步给出的就是 value function 的理论解。换言之，固定价格 Bellman 模块不只是“有解”，而且该解在给定的 benchmark primitives 与 regularity assumptions 下是唯一的。

\paragraph{第五步：用 value iteration 构造性地计算该 fixed point。}
Banach 定理不仅给出存在唯一性，也同时给出构造性算法。对任意有界初值 $V^{(0)}$，反复迭代
\[
V^{(k+1)}=\mathcal T_MV^{(k)}
\]
就会收敛到唯一 fixed point $V^*$。在离散网格上，这个过程分成三步：
\begin{enumerate}
    \item 先用转移矩阵计算 continuation values；
    \item 再逐点计算四个动作的 choice-specific values；
    \item 最后逐点取最大值，得到新的 value-function iterate。
\end{enumerate}
如果记
\[
\Delta_k:=\|V^{(k+1)}-V^{(k)}\|,
\]
那么压缩性质给出
\[
\Delta_k\le \beta\Delta_{k-1}\le \beta^k\Delta_0.
\]
所以 successive differences 以几何速度收缩，迭代序列是 Cauchy 的，并且收敛到唯一的 $V^*$。这一步完成了从理论 fixed point 到可执行求解算法的转换，因此它是后续数值实现与实证接口的直接入口。

\paragraph{第六步：在 value function 求解之后恢复 policy objects。}
只有当 value iteration 已经收敛，或者在理论上已经得到唯一 fixed point $V^*$ 之后，才应该去计算
\[
\Gamma_s(z;M)=\arg\max_{a\in\mathcal A}\mathcal W_{sa}(z;M;V^*)
\]
以及
\[
\mathcal Z_{sj}(M)=\{z\in Z:j\in\Gamma_s(z;M)\}.
\]
顺序不能反过来，因为 policy objects 是由 solved value function 派生出来的，而不是先于 value function 独立存在的原始对象。更进一步地说，Phase 4 所需要的组织转移矩阵与 stationary distribution block，正是以这些 choice regions 作为输入来构造的。因此，Phase 3 的求解终点不是停留在“值函数存在”，而是要把唯一的 $V^*$ 转化为可测的策略与区域对象，使下一阶段的 transition aggregation 和 stationary distribution analysis 具有明确输入。

\paragraph{形式化总结。}
在本 benchmark 中，Bellman/value-function 的求解顺序可以概括为：
\begin{align*}
\text{先写 }W_{sj}
&\Longrightarrow
\text{再写 }V_s=\max_j W_{sj}\\
&\Longrightarrow
\text{把问题改写成 }V=\mathcal T_MV\\
&\Longrightarrow
\text{用 Banach 定理得到唯一 fixed point}\\
&\Longrightarrow
\text{用 value iteration 构造性地求出 }V^*\\
&\Longrightarrow
\text{最后恢复 policy correspondence 与 choice regions。}
\end{align*}
而这些 policy correspondence 与 choice regions 又构成了 Phase 4 中 transition matrix、stationary distribution 与后续 equilibrium objects 的直接输入。

\subsection{Bellman 模块数值求解（完整步骤）}

固定点定理告诉我们值函数存在且唯一，但一年级 PhD 学生仍然需要看到这个固定点究竟怎样被算出来。由于 $\mathcal T_M$ 是压缩映射，构造性算法就是 value iteration。

\paragraph{步骤 N1：选择状态的数值表示并初始化。}
选取网格点或求积节点
\[
z_1,\dots,z_N\in Z
\]
来逼近紧状态空间。对每个活跃类型 $j\in\{A,B,C\}$，构造一个 $N\times N$ 的 row-stochastic 矩阵 $\mathbf Q_j=(q_{j,nm})$，满足
\[
q_{j,nm}\ge 0,
\qquad
\sum_{m=1}^N q_{j,nm}=1
\qquad\text{对每个 }n.
\]
这个性质就是“$Q_j(\cdot\mid z,M)$ 是概率测度”的离散版本。

随后选择任意一个有界初值
\[
V^{(0)}=(V_A^{(0)},V_B^{(0)},V_C^{(0)}).
\]
一个方便的 benchmark 初值是
\[
V_s^{(0)}(z_n)=0
\qquad\text{对所有 }s\in\{A,B,C\},\ n=1,\dots,N.
\]
任何其他有界初值也都可以，因为压缩映射保证从任意起点出发都收敛。

\paragraph{步骤 N2：计算 continuation value。}
对每个 $j\in\{A,B,C\}$ 和每个网格节点 $z_n$，计算
\[
\mathcal C_j^{(k)}(z_n):=\sum_{m=1}^N q_{j,nm}V_j^{(k)}(z_m).
\]
这是积分
\[
\int_Z V_j^{(k)}(z')Q_j(dz'\mid z_n,M)
\]
的离散近似。在这一步里，还没有做最大化；我们只是在当前 iterate 下评估 continuation value。

\paragraph{步骤 N3：逐点计算 choice-specific values。}
对每个继承类型 $s$ 和每个节点 $z_n$，计算
\begin{align*}
W_{sA}^{(k)}(z_n)
&=
\Pi_A(z_n;w,M)-\kappa_{sA}(M)+\beta\mathcal C_A^{(k)}(z_n),\\
W_{sB}^{(k)}(z_n)
&=
\Pi_B(z_n;p_m,M)-\kappa_{sB}(M)+\beta\mathcal C_B^{(k)}(z_n),\\
W_{sC}^{(k)}(z_n)
&=
\Pi_C(z_n;p_m,w,M)-\kappa_{sC}(M)+\beta\mathcal C_C^{(k)}(z_n),\\
W_{s,\mathrm{exit}}^{(k)}(z_n)&=0.
\end{align*}
这里没有任何隐藏步骤。每一个数值都等于“当前优化收益减去当前切换成本，再加上折现后的期望 continuation value”。

\paragraph{步骤 N4：做一次 Bellman 更新。}
对每个继承类型 $s$ 和每个节点 $z_n$，令
\[
V_s^{(k+1)}(z_n)
=
\max\Bigl\{
W_{sA}^{(k)}(z_n),
W_{sB}^{(k)}(z_n),
W_{sC}^{(k)}(z_n),
0
\Bigr\}.
\]
这一步就是把 Bellman 算子在离散状态空间上应用一次。

\paragraph{步骤 N5：定义计算时真正可观测的误差量。}
每做完一次更新，就计算 successive-difference residual
\[
\Delta_k:=\|V^{(k+1)}-V^{(k)}\|.
\]
数值上，它就是
\[
\Delta_k=
\max_{s\in\{A,B,C\}}
\max_{n=1,\dots,N}
\left|V_s^{(k+1)}(z_n)-V_s^{(k)}(z_n)\right|.
\]
这个量在算法内部是可以直接观察到的。

\paragraph{步骤 N6：为什么 successive differences 以几何速度收缩。}
由于 $V^{(k+1)}=\mathcal T_MV^{(k)}$ 且 $V^{(k)}=\mathcal T_MV^{(k-1)}$，有
\[
\Delta_k
=
\|\mathcal T_MV^{(k)}-\mathcal T_MV^{(k-1)}\|
\le \beta\|V^{(k)}-V^{(k-1)}\|
=
\beta\Delta_{k-1}.
\]
递推下去，得到
\[
\Delta_k\le \beta^k\Delta_0.
\]
因此，相邻两次 iterate 之间的差距以几何速度收缩。

\paragraph{步骤 N7：为什么值函数序列收敛。}
取任意 $m>n$，由三角不等式可得
\begin{align*}
\|V^{(m)}-V^{(n)}\|
&\le
\sum_{r=n}^{m-1}\|V^{(r+1)}-V^{(r)}\|\\
&=
\sum_{r=n}^{m-1}\Delta_r\\
&\le
\sum_{r=n}^{m-1}\beta^r\Delta_0\\
&\le
\frac{\beta^n}{1-\beta}\Delta_0.
\end{align*}
因为 $\beta\in(0,1)$，当 $n\to\infty$ 时，右边趋于零。因此 $\{V^{(k)}\}$ 是 Cauchy 序列。由于值函数空间完备，该序列收敛到某个极限。再由 Banach 定理，这个极限必然就是唯一的 Bellman 固定点 $V^*$。

\paragraph{步骤 N8：a-posteriori error bound 与停机准则。}
上面的推导还给出一个可计算的误差上界。对任意 $k$，
\begin{align*}
\|V^{(k)}-V^*\|
&\le
\sum_{r=k}^{\infty}\|V^{(r+1)}-V^{(r)}\|\\
&=
\sum_{r=k}^{\infty}\Delta_r\\
&\le
\sum_{m=0}^{\infty}\beta^m\Delta_k\\
&=
\frac{1}{1-\beta}\Delta_k.
\end{align*}
因此，一个有效的停机准则是
\[
\frac{\Delta_k}{1-\beta}<\varepsilon_{\mathrm{fp}},
\]
其中 $\varepsilon_{\mathrm{fp}}>0$ 是预先设定的 fixed-point tolerance。

\paragraph{步骤 N9：值函数收敛后再恢复策略对象。}
只有在 value iteration 已经收敛之后，才用收敛后的 $V^*$ 去恢复 benchmark 的策略对象：
\[
\Gamma_s(z_n;M)=\arg\max_{a\in\mathcal A}W_{sa}^{(*)}(z_n),
\qquad
\mathcal Z_{sj}(M)=\{z_n:j\in\Gamma_s(z_n;M)\}.
\]
如果后续计算必须使用 deterministic policy，那么也应该在 benchmark correspondence 已经算出之后，再额外施加明确声明的 tie-breaking 规则。

\subsection{What did we just do, and why does it matter?}

我们从 Phase 2 的优化后静态收益对象出发，把它们嵌入了一个动态组织选择问题。这里最关键的概念区分是：组织选择是外层动态选择，类型专属的经营控制是内层期内选择，而这些内层问题已经在 Phase 2 中被压缩为 $\Pi_A$、$\Pi_B$、$\Pi_C$。所以 Bellman 系统比较的是 $W_{sA}$、$W_{sB}$、$W_{sC}$ 和 exit，而不是重新把静态控制问题写一遍。

接着，我们把三条 Bellman 方程改写成定义在完备度量空间上的一个算子方程。这一步很重要，因为它把一个相互递归的方程组转换成了一个固定点问题。只要 Bellman 算子被定义清楚，存在性问题就变成“这个算子有没有固定点”，唯一性问题就变成“这个算子会不会有多个固定点”。

压缩性质是整个论证的核心。它告诉我们，Bellman 算子会把任意两个候选值函数三元组之间的距离按比例 $\beta$ 缩小。这个结论同时带来两个结果。第一，固定点只能有一个。第二，从任意有界初值出发的反复 Bellman 更新都会收敛到这个固定点。也就是说，同一个数学论证同时给出了 well-posedness 和可执行的求解算法。

值函数唯一性之所以重要，是因为后面几乎所有对象都依赖于它。argmax correspondence 依赖于 $V^*$。choice regions 依赖于 argmax correspondence。稳态分布模块里的组织转移矩阵依赖于 choice regions。entry value 也依赖于 value functions。如果 $V$ 不是唯一的，那么这些后续对象也会变得不唯一。

\subsection{End-of-Section Recap}

\paragraph{What has been established mathematically.}
对固定的 $(w,p_m,M)$，benchmark 的动态组织选择问题在 $\mathcal V=\mathcal C_b(Z)^3$ 上定义了一个 Bellman 算子。在给定的正则条件下，这个算子是压缩映射，因此 Bellman 系统存在唯一的值函数三元组。benchmark 的 argmax correspondence 与 choice regions 都有良好定义，而且 value iteration 以显式可计算的误差上界收敛到唯一固定点。

\paragraph{What assumptions were used.}
我们使用了精确的 benchmark payoff 对象 $\Pi_A(z;w,M)$、$\Pi_B(z;p_m,M)$、$\Pi_C(z;p_m,w,M)$，切换成本矩阵 $\kappa_{sj}(M)$，continuation kernels $Q_j(\cdot\mid z,M)$，以及 $\beta\in(0,1)$ 的折现。为完成证明，还额外使用了：状态空间 $Z$ 是非空紧致 metric space；payoff 对象有界连续；转移核满足 Feller 连续性；切换成本有限。

\paragraph{What remains to be derived next.}
下一阶段需要利用这里得到的唯一固定价格值函数以及 benchmark choice regions，进一步构造内生转移矩阵、各组织类型的稳态分布、entry block，以及共同决定 $(w,p_m)$ 与长期配置的 stationary equilibrium 条件。


